\begin{proposition}[Pointwise cusp limits and the fixed double curve]\hypertarget{tag-02PB}{}\label{mapsI:src:A.1.2}
\StatementRecord{mapsI:src:A.1.2}{proved}{02PB}{\SeriesRef{mapsI:quotient-eq-4}}

The literal character has no continuous extension at either
triple point or at generic points of two of the three double
curves.  It has the constant limit $0$ on the smooth locus of
$W$ and on the third double curve away from the triple points.
In the chart $q=z_0z_1z_2$, $x_1=z_1$, $x_2=z_2$, the extendible
double curve is $z_0=z_1=0$, $z_2\ne0$.  The fixed curve of the
$\mathbb C^*$-action in direction $d=e_2$ is instead
$z_0=z_2=0$, $z_1\ne0$, and is one of the nonextendible curves.
Consequently the extendible double curve and the fixed double curve are distinct.
\end{proposition}



\begin{proof}
\begin{enumerate}[label=\arabic*.,leftmargin=*]
\item Put $q=e^{-2\pi T}$ on a positive radial ray and use the exact
cusp periods
\[
 Z_T=iTB_0+C(q),\qquad
 B_0=\begin{pmatrix}0&1\\-1&0\end{pmatrix},
\]
where $C$ is bounded and holomorphic at $0$.  For a fixed
$a\in\mathbb R^2$ the representative
$x_T(a)=\exp(2\pi iZ_Ta)$ has character $a_1\bmod\mathbb Z$ and
logarithmic position
\begin{equation}
 \frac{\log|x_T(a)|}{\log|q|}
 =B_0a+T^{-1}\operatorname{Im}C(q)a.
 \hypertarget{tag-02P0}{}\label{mapsI:quotient-eq-3}
\end{equation}
The choices $a=(-1/4,1/4)$ and $a'=(-1/2,1/4)$ have logarithmic
positions tending to $(1/4,1/4)$ and $(1/4,1/2)$.
Both are interior to the triangle with vertices $0,e_1,e_2$,
so all three chart coordinates tend to $0$.  Their character
values are $3/4$ and $1/2$.  This proves nonextension at the
first triple point.  The other triangle type has vertices
$e_1,e_2,e_1+e_2$.  Its coordinates can be taken as
$x_1x_2/q,q/x_2,q/x_1$.  Positions $(3/4,3/4)$ and
$(3/4,1/2)$ lie in its interior, give distinct values of $-y_2$
modulo one, and converge to its origin.  Translation by the
lattice preserves triangle type; these are precisely the two
triple points of $W$.

\item For arbitrary convergent paths in the first chart, inversion
of the exact logarithmic period matrix gives
\begin{equation}
 \theta_\gamma=-\frac{\log|z_2|}{\log|q|}+o(1)\pmod{\mathbb Z}.
 \hypertarget{tag-02P1}{}\label{mapsI:quotient-eq-4}
\end{equation}
Here the error tends to zero whenever the logarithmic ratios
stay bounded; they do stay bounded near a fixed point of a
normal-crossing stratum because $q=z_0z_1z_2$ and the three
coordinates are bounded.  If $z_0,z_1\to0$ and $z_2\to c\ne0$,
the right side tends to zero.  If $z_0,z_2\to0$ and $z_1=c\ne0$,
the paths $z_2=q^a,z_0=q^{1-a}/c$, $0<a<1$, have limits
$-a\pmod1$ at the same double point.  The same construction
with $z_0=c$ proves nonextension on $z_1=z_2=0$.
At a smooth point of a component the ratio in \eqref{mapsI:quotient-eq-4} tends
to $0$ or $1$, giving zero modulo one.  Integer translates
of these charts give the stated assertion on all of $W$.

\item The torus action in direction $e_2=d$ is
\[
 \lambda\cdot(z_0,z_1,z_2)
       =(\lambda^{-1}z_0,z_1,\lambda z_2).
\]
Its fixed curve is $z_0=z_2=0$, while the character extends along
$z_0=z_1=0$ by the preceding limit formula. The displayed action and
limit formula give these two loci in the same toric chart.
\end{enumerate}
\end{proof}